\documentclass[11pt]{article}
\usepackage[margin=1in]{geometry}
\usepackage{amsmath,amssymb,amsthm,hyperref}
\newtheorem{theorem}{Theorem}
\newtheorem{lemma}[theorem]{Lemma}
\newtheorem{proposition}[theorem]{Proposition}
\newcommand{\R}{\mathcal R}
\newcommand{\lcm}{\operatorname{lcm}}
\title{Asymptotically full minimal signed unit-fraction identities\\
\large An alternative proof for Erd\H{o}s Problem 319}
\author{\begin{tabular}{cc}
Jensen Kohlmeyer & Liam Kruer\\
\texttt{jenw1n} & \texttt{lkruer}\\[4pt]
Will Howe & Grant Huffman\\
\texttt{H0W4RD-dev} & \texttt{Ghuffy1}
\end{tabular}}
\date{16 October 2026}
\begin{document}
\maketitle
\begin{abstract}
Let $f(N)$ be the largest size of a subset of $\{1,\ldots,N\}$ admitting
a signed zero sum of its reciprocals with no proper nonempty signed
zero subsum. We prove $f(N)=N-o(N)$, using the subset-sum theorem of
Liu and Sawhney. The construction joins finitely many large-prime groups
by one CRT-chosen unit fraction. The resulting positive sum can be
arbitrarily large, although the density cost of excluding those primes
is arbitrarily small. A dense coprime-denominator expansion of that
sum provides the opposite side and preserves minimality.
This is a complete candidate proof, self-checked but not independently
reviewed. Earlier public work already claims this asymptotic and the
protected-prime construction used here. This version gives a route through
Liu--Sawhney and an accompanying Lean formalization; no first-discovery or
formalization-priority claim is made. Substantial AI assistance was used.
\end{abstract}

\section{Statement}
For a finite set $B\subset\mathbb N$, write $\R(B)=\sum_{n\in B}1/n$.
Call $A\subseteq\{1,\ldots,N\}$ primitive if there are signs
$\delta_n\in\{-1,1\}$ satisfying
\[
 \sum_{n\in A}\frac{\delta_n}{n}=0,
 \qquad
 \sum_{n\in C}\frac{\delta_n}{n}\ne0
 \quad\text{for every }\varnothing\ne C\subsetneq A.
\]
Let $f(N)$ be the maximum of $|A|$ over these sets.

\begin{theorem}\label{main}
For every $\varepsilon>0$ and every sufficiently large $N$, there is a
primitive $A\subseteq\{1,\ldots,N\}$ with $|A|\ge(1-\varepsilon)N$.
Consequently
\[
 \boxed{\displaystyle\lim_{N\to\infty}\frac{f(N)}N=1.}
\]
\end{theorem}

The lower bound $(1-1/e+o(1))N$ follows from a dense expansion of $1$
and the singleton $\{1\}$ \cite{problem}. Earlier public candidate
proofs of the density-one asymptotic are recorded in
\cite{g01234,peilinliu}; their existence precludes treating the present
result as an uncontested new solution.
Our positive side is a fixed, possibly enormous, finite set once
$\varepsilon$ is chosen. Only the opposite side grows with $N$.
The result is asymptotic and supplies no practical threshold for $N$.

\section{A finite positive set with only one coprime subsum}
For a squarefree integer $D$, call a rational number $D$-integral if its
reduced denominator is coprime to $D$. This includes $0$ and all integers.
The $D$-integral rationals form an additive subgroup of $\mathbb Q$.

\begin{lemma}\label{crt}
Let $K\ge1$, $L=\lcm(1,\ldots,K)$, and $B=LH_K$, where
$H_K=\sum_{m=1}^K1/m$.
Let $\mathcal P$ be a finite set of at least two primes, each greater
than $B$, and put $D=\prod_{p\in\mathcal P}p$.
There exists a positive integer $t$ coprime to $D$ such that the set
\[
 E=\{pm:p\in\mathcal P,\ 1\le m\le K\}\ \cup\ \{Dt\}
\]
has the following properties:
\begin{enumerate}
\item All its denominators are distinct, and all have a prime factor in
$\mathcal P$.
\item $r:=\R(E)$ is $D$-integral and
\[
 r=H_K\sum_{p\in\mathcal P}\frac1p+\frac1{Dt}.
\]
\item The only subsets $C\subseteq E$ for which $\R(C)$ is
$D$-integral are $\varnothing$ and $E$.
\end{enumerate}
\end{lemma}
\begin{proof}
Since $p>B\ge K$, both $L$ and $B$ are invertible modulo every
$p\in\mathcal P$. Choose $t$ by the Chinese remainder theorem so that
\begin{equation}\label{crt-cong}
 tB(D/p)+L\equiv0\pmod p
 \qquad(p\in\mathcal P).
\end{equation}
The prescribed residue $-L\bigl(B(D/p)\bigr)^{-1}$ is nonzero modulo
$p$, so $\gcd(t,D)=1$. Take its positive representative, or add a
multiple of $D$.

The groups $\{pm:m\le K\}$ are disjoint: a denominator in the $p$-group
has no other prime factor in $\mathcal P$. The connector $Dt$ has all
these prime factors, so is distinct from every group denominator.

For a subset $C\subseteq E$, let
\[
 J_p=\{m\le K:pm\in C\},\qquad
 b_p=\sum_{m\in J_p}L/m,
\]
and let $h\in\{0,1\}$ indicate whether $Dt\in C$.
Every $b_p$ is an integer in $[0,B]$, with $b_p=0$ exactly for an empty
$J_p$, and $b_p=B$ exactly for a full $J_p$.
Multiplying $\R(C)$ by the common denominator $DtL$ gives the integer
\begin{equation}\label{numerator}
 U_C=t\sum_{p\in\mathcal P}(D/p)b_p+hL.
\end{equation}
Each $p\in\mathcal P$ occurs to exponent exactly one in $DtL$.
Thus $\R(C)$ is $D$-integral if and only if every such $p$ divides $U_C$.

If $h=0$, reduction modulo $p$ gives $b_p\equiv0\pmod p$.
As $0\le b_p\le B<p$, this forces $b_p=0$ for every $p$, hence $C$ empty.
If $h=1$, subtracting \eqref{crt-cong} from the condition $p\mid U_C$
gives $b_p\equiv B\pmod p$. The same bound forces $b_p=B$ for every $p$,
hence $C=E$.
Finally, \eqref{crt-cong} directly makes $U_E$ divisible by every $p$,
so $r$ is $D$-integral. The formula for $r$ is immediate.
\end{proof}

\begin{lemma}\label{amplification}
For every $0<\delta<1/4$ and $R>0$, there are a squarefree $D$ and a
finite set $E$ as in Lemma~\ref{crt} such that
\[
 \rho:=\frac{\varphi(D)}D\ge1-2\delta,
 \qquad r=\R(E)\ge R.
\]
\end{lemma}
\begin{proof}
Choose $K$ so that $H_K\ge R/\delta$, and then choose primes
$p>\max\{LH_K,1/\delta\}$ until
\[
 \delta\le\sum_{p\in\mathcal P}\frac1p<2\delta.
\]
Such a finite choice exists because the sum of reciprocal primes
diverges. Each summand is less than $\delta$, so at least two primes
are chosen. The finite product inequality gives
\[
 \rho=\prod_{p\in\mathcal P}(1-1/p)
 \ge1-\sum_{p\in\mathcal P}1/p>1-2\delta.
\]
Lemma~\ref{crt} gives $r\ge H_K\delta\ge R$.
\end{proof}

\section{Dense expansions avoiding a fixed set of primes}
The only deep input is Proposition 5.2 of Liu--Sawhney \cite{LS}.
For transparency, its hypotheses are recorded in the appendix. The
following application also explains why fixed prime factors in the
target denominator cause no problem.

\begin{lemma}\label{dense}
Fix a squarefree $D$ and a positive $D$-integral rational $r$.
Put $\rho=\varphi(D)/D$.
For every $0<\kappa<\min(r/2,1/2)$, set $c=\exp(-(r+\kappa)/\rho)$.
For all sufficiently large $N$, there is a set
\[
 F_N\subseteq\{n\in[cN,N]:\gcd(n,D)=1\}
\]
with $\R(F_N)=r$ and
\begin{equation}\label{density}
 |F_N|\ge\bigl(\rho(1-c)-\kappa+o(1)\bigr)N.
\end{equation}
In particular,
\[
 \liminf_{N\to\infty}\frac{\max|F_N|}N
 \ge\rho(1-e^{-r/\rho}),
\]
where the maximum allows all coprime denominators at most $N$.
\end{lemma}
\begin{proof}
Let $W_N=\{n\in[cN,N]:\gcd(n,D)=1\}$.
Since $D,c$ are fixed, elementary counting in residue classes gives
\begin{equation}\label{original-mass}
 |W_N|=\rho(1-c)N+O_D(1),\qquad
 \R(W_N)=\rho\log(1/c)+o(1)=r+\kappa+o(1).
\end{equation}
Set
\[
 M_0=\min(c/2,10^{-5})N,\quad
 u=(\log N)^{4/5},\quad K_0=M_0e^{-u},\quad
 S=Ne^{-6u},\quad \eta=(\log\log N)^{-2}.
\]
The endpoint can be rounded to an integer. Remove from $W_N$ all
numbers that are not $S$-smooth and all numbers with
$\Omega(n)>5\log\log N$. Lemmas 2.2--2.3 of \cite{LS} show that this
removes $O(N/(\log N)^{1/5})+O(N/(\log N)^3)=o(N)$ elements.
Because these numbers are at least $cN$, their reciprocal mass is $o(1)$.

Next, while a prime power $q$ occurring in the current set satisfies
$q\R(W_q)<\eta$, remove all its multiples. Each prime power is used
at most once, so the total reciprocal mass removed is at most
\[
 \eta\sum_{q\le N}\frac1q=O(1/\log\log N)=o(1).
\]
Call the final set $W^*_N$, and put $Q=\lcm(W^*_N)$.
We have
\begin{equation}\label{cleaned}
 \R(W^*_N)=r+\kappa+o(1),\qquad
 q\R((W^*_N)_q)\ge\eta
\end{equation}
for every prime power occurring in $W^*_N$.

Write $r=a/b$ in lowest terms. For each prime power $q\mid b$,
$\gcd(q,D)=1$. Again by residue-class counting,
\[
 \R((W_N)_q)=\frac{\rho}{q}\log(1/c)+o(1)>0.
\]
Here $q$ is fixed. Since the entire cleanup removes only $o(1)$
reciprocal mass, some multiple of $q$ remains in $W^*_N$.
Therefore $b\mid Q$. In particular,
\[
 x=Q\bigl(\R(W^*_N)-r\bigr)
\]
is an integer, positive for large $N$, and $x/Q=\kappa+o(1)$.

Apply Proposition 5.2 with $d=1/5$, $e=1/100$, and the displayed
$M_0,K_0,S,\eta$. All its size conditions hold for large $N$.
For the last condition, the second term in $G$ gives
\[
 G\gg_{c}(\log N)^{3/5}(\log\log N)^{-9},
\]
and hence
\[
 \frac{G^2}{(\log N)^{2e}
  (\log(N/M_0)+(\log N)^{4/5})}
 \gg_c\frac{(\log N)^{19/50}}{(\log\log N)^{18}}
 \longrightarrow\infty.
\]
The bounds on $S$ hold because its factor $e^{-6u}$ dominates the
$e^{-2u}$ and polylogarithmic factors in those bounds;
$S\ge N^{.99}$ eventually. The target condition follows from
\eqref{cleaned} and $\kappa>0$:
\[
 (1+1/\log N)x/Q\le\R(W^*_N)\le(\log N)x/Q.
\]
The proposition supplies $G_N\subseteq W^*_N$ with reciprocal sum
$x/Q=\R(W^*_N)-r$. Take $F_N=W^*_N\setminus G_N$, so that
$\R(F_N)=r$. This applies the proposition to a small complementary
target; its sampling proportion is less than $1/2$ for large $N$.
Finally every term in $W_N\setminus F_N$ is at most $N$, so
\[
 |W_N\setminus F_N|
 \le N\R(W_N\setminus F_N)
 =N(\kappa+o(1)).
\]
Together with \eqref{original-mass} this proves \eqref{density}.
Taking the lower limit in $N$ and then letting $\kappa\downarrow0$
gives the final assertion. No uniformity as $D,r,\kappa$ vary is required.
\end{proof}

\section{Primitivity and completion of the asymptotic}
\begin{proof}[Proof of Theorem~\ref{main}]
Fix $0<\delta<1/4$ and use Lemma~\ref{amplification} with
$R=\log(1/\delta)$. This gives fixed $D,E,r$ and
$\rho\ge1-2\delta$, $r\ge\log(1/\delta)$.
By Lemma~\ref{dense}, choose $F_N$ with $\R(F_N)=r$ and all
denominators coprime to $D$. For large $N$, every element of the fixed
set $E$ is at most $N$. Also $E\cap F_N=\varnothing$, since every
element of $E$ has a prime factor dividing $D$.

Give $E$ positive signs and $F_N$ negative signs. The total is zero.
Suppose $C\subseteq E\cup F_N$ is a signed zero subset, and put
$E'=C\cap E$, $F'=C\cap F_N$. Then
\[
 \R(E')=\R(F').
\]
The right side is $D$-integral. Lemma~\ref{crt} therefore forces
$E'=\varnothing$ or $E'=E$.
In the first case positivity forces $F'=\varnothing$.
In the second, $\R(F')=r=\R(F_N)$ and positivity forces $F'=F_N$.
Thus the signed identity is primitive.

Its cardinality and Lemma~\ref{dense} give
\[
 \liminf_{N\to\infty}\frac{f(N)}N
 \ge\rho(1-e^{-r/\rho})
 \ge(1-2\delta)(1-\delta)\ge1-3\delta.
\]
The middle inequality uses $\rho\le1$ and
$r/\rho\ge r\ge\log(1/\delta)$.
Let $\delta\downarrow0$, and combine this with $f(N)\le N$.
This proves the asserted limit and the quantified first statement.
\end{proof}

For a small concrete example, $K=2$ and $\mathcal P=\{5,7,11\}$ give
$D=385$, $t=73$, and the primitive identity
\[
 \frac15+\frac1{10}+\frac17+\frac1{14}+\frac1{11}+\frac1{22}
 +\frac1{28105}
 =\frac{95}{146}
 =\frac12+\frac18+\frac1{39}+\frac1{22776}.
\]
The four denominators on the right are coprime to $385$.
This illustrates minimality, not the asymptotic density estimate.

\section{Verification and attribution}
The protected-prime construction, the amplification of its reciprocal
mass at small prime-density cost, and the final primitivity argument
already appear in the public manuscript of \texttt{g-01234}, dated
16 July 2026 \cite{g01234}. A later submission by
\texttt{peilinliu66-dev}, opened on 3 October 2026, claims a complete
Lean proof using a thinned-prime inverse-sum completion route
\cite{peilinliu}. That submission was pending when checked on
16 October 2026; we have not independently reproduced it. These earlier
claims must be considered in any assessment of priority or prize eligibility.
The present dense-expansion lemma instead applies Liu--Sawhney's
Proposition 5.2. No novelty claim is made for the protected-prime idea.
The Fourier analysis is due to Liu--Sawhney and their predecessors.
The attached exact audit checks small global CRT examples exhaustively
and checks a larger example locally at every selected prime. These
finite checks do not replace the proof of the asymptotic or certify
the imported analytic theorem.

The accompanying file \texttt{Erdos319.lean} proves the limit and explicit
fixed-sign witnesses without additional mathematical hypotheses. Its
analytic estimates include ports of existing proofs from
\texttt{plby/lean-proofs}, commit
\texttt{8822f7ddef30fadbd92e1c6ab4ed897af356af5e}; those background proofs
retain their separate attribution. The checked environment is Lean
4.35.0-rc2 with Mathlib commit
\texttt{065356127b1dc0016f66b7283ce0ce2c4055aa55}.
The main theorem is
\begin{center}
\texttt{Erdos319Formal.erdos\_problem\_319}.
\end{center}
The original 1980 statement was checked at printed p.~43 \cite{EG}:
the same signs are used in the whole sum and in every subsum.
The conclusion determines the leading asymptotic, not the lower-order
deficit or the exact maximum for each finite $N$.

The listed team prepared this AI-assisted proof presentation and
formalization for review. Team credit does not assert that each member
manually wrote every proof step. No independent human review, prize
acceptance, or entitlement to payment is asserted.

\appendix
\section{The precise Liu--Sawhney input}
Here is Proposition 5.2 of \cite{LS}, restated with renamed parameters.
For $W\subseteq[1,N]$, let $W_q=\{n\in W:q\mid n\}$ and let $Q_W=\lcm(W)$.
There is an absolute $C$ such that, for fixed $d\in(0,1)$,
$e\in(0,1/10)$ and sufficiently large $N$, the following holds.
Suppose $\eta\ge1/\log N$ and
\[
 N^{.99}\le S\le K_0\le M_0\le N/10^4,\qquad
 K_0\le M_0\exp(-(\log N)^{1-d}),
\]
\[
 S\le\min\left\{\frac{M_0^2}{CN},
 \frac{\eta M_0K_0^2}{N^2(\log N)^3}\right\}.
\]
Define
\[
 G=\max\left\{
 \frac{\eta}{(\log N)^d(\log\log N)^3},
 \frac{\eta^2(\log N)^{1-2d}}
 {\log(N/M_0)^2(\log\log N)^5}\right\}.
\]
Assume
\[
 \frac{G^2}{(\log N)^{2e}
  (\log(N/M_0)+(\log N)^{1-d})}\ge C.
\]
Let $W\subseteq[M_0,N]$ satisfy: every prime power dividing an element
of $W$ is at most $S$; $\Omega(n)\le5\log\log N$ for $n\in W$; and
$q\R(W_q)\ge\eta$ for every prime power occurring in $W$.
For every positive integer $x$ with
\[
 (1+1/\log N)x/Q_W\le\R(W)\le(\log N)x/Q_W,
\]
there is $W'\subseteq W$ such that $\R(W')=x/Q_W$.

In the dense-expansion lemma all parameters other than $N$ are fixed
or explicitly specified functions of $N$. In particular, the proof
does not use the inconsistent displayed sign of $K$ in the arXiv
version's proof of Theorem 1.3; it uses the proposition's stated
$K_0\le M_0\exp(-(\log N)^{1-d})$ condition directly.

\begin{thebibliography}{9}
\bibitem{LS} Y. P. Liu and M. Sawhney,
\emph{On further questions regarding unit fractions},
\href{https://arxiv.org/html/2404.07113v1}{arXiv:2404.07113v1} (2024),
Proposition 5.2 and Lemmas 2.2--2.3.
International Mathematics Research Notices (2026), rnaf382;
\href{https://doi.org/10.1093/imrn/rnaf382}{published article}.
\bibitem{EG} P. Erd\H{o}s and R. L. Graham,
\emph{Old and new problems and results in combinatorial number theory}
(1980), p.~43;
\href{https://mathweb.ucsd.edu/~ronspubs/80_11_number_theory.pdf}{original scan}.
\bibitem{problem} T. F. Bloom,
\href{https://www.erdosproblems.com/319}{Erd\H{o}s Problem 319}.
\bibitem{g01234} \texttt{g-01234},
\emph{A Candidate Proof That $c(N)/N\to1$ in Erd\H{o}s Problem 319},
circulation draft, 16 July 2026;
\href{https://github.com/g-01234/erdos-319-density-one/blob/f55a25a29dcb17193918d27506be5e7a5b5d7d40/docs/proof.md}{pinned manuscript}.
\bibitem{peilinliu} \texttt{peilinliu66-dev},
\emph{JSP-000265: complete Erd\H{o}s 319 fixed-sign density-one Lean proof},
submission opened 3 October 2026;
\href{https://github.com/TheJustinSunPrize/awards/pull/4669}{awards PR 4669};
\href{https://github.com/peilinliu66-dev/jsp265-erdos319-density-one/tree/72545d3e68e8881e2f89f4744067583aaa61fc9e}{selected proof revision}.
\end{thebibliography}
\end{document}
